\subsection{Direct image and inverse image of a sheaf}

Let A and B be topological spaces and \(u: A \rightarrow B\) a continuous map.

Let \(\mathcal{F} = (\mathcal{F}(\mathrm{U}), f_{\mathrm{UV}})\) be a presheaf on A. We define a presheaf \(\mathcal{F}'\) on B as follows: for every open set U of B, set \(\mathcal{F}'(\mathrm{U}) = \mathcal{F}(u^{-1}(\mathrm{U}))\) and for every pair (U, V) of open sets of B such that U \(\subset\) V, set \(f_{\mathrm{UV}}' = f_{u^{-1}(\mathrm{U})u^{-1}(\mathrm{V})}\) . Then \((\mathcal{F}'(\mathrm{U}), f_{\mathrm{UV}}')\) is a presheaf on B. We denote it by \(u_{*}(\mathcal{F})\) and call it the direct image presheaf of the presheaf \(\mathcal{F}\) under the map u.

If \((\mathrm{U}_{i})_{i\in\mathrm{I}}\) is a family of open sets of B, we have \(u^{-1}(\bigcup_{i\in\mathrm{I}}\mathrm{U}_{i})=\bigcup_{i\in\mathrm{I}}u^{-1}(\mathrm{U}_{i})\) and \(u^{-1}(\bigcap_{i\in\mathrm{I}}\mathrm{U}_{i})=\bigcap_{i\in\mathrm{I}}u^{-1}(\mathrm{U}_{i})\) (E, II, p.~25, prop. 3 and 4). It follows immediately that if F has property \((\mathrm{F}_{1})\) (resp. \((\mathrm{F}_{2})\) ) of sheaves (I, p.~43), then so does \(u_{*}(\mathcal{F})\) . Consequently, the direct image of a sheaf is a sheaf.

Let \(\mathcal{F}_1\) and \(\mathcal{F}_2\) be presheaves on A and let \(\varphi\colon\mathcal{F}_1\to\mathcal{F}_2\) be a morphism of presheaves. There then exists a unique morphism of presheaves \(u_*\varphi\colon u_*\mathcal{F}_1\to u_*\mathcal{F}_2\) such that for every open set U of B, the map \((u_*\varphi)(U)\colon(u_*\mathcal{F}_1)(U)\to(u_*\mathcal{F}_2)(U)\) is the map \(\varphi(u^{-1}(U))\colon\mathcal{F}_{1}(u^{-1}(U))\to\mathcal{F}_{2}(u^{-1}(U)).\) If \(\mathcal{F}_{3}\) is a presheaf on A and if \(\psi\colon\mathcal{F}_{2}\to\mathcal{F}_{3}\) is a morphism of presheaves, we have \(u_{*}(\psi\circ\varphi)=u_{*}(\psi)\circ u_{*}(\varphi)\) .

Let C be a topological space and let \(v: B \to C\) be a continuous map. If \(\mathcal{F}\) is a presheaf on A, the presheaves \(v_{*}(u_{*}(\mathcal{F}))\) and \((v \circ u)_{*}(\mathcal{F})\) coincide. If \(\varphi: \mathcal{F}_{1} \to \mathcal{F}_{2}\) is a morphism of presheaves on A, we have the equality \(v_{*}(u_{*}(\varphi)) = (v \circ u)_{*}(\varphi)\) .

Example 1. --- Let B be a topological space, A a subspace of B, and \(\mathcal{F} = (\mathcal{F}(U), f_{\mathrm{UV}})\) a presheaf on A. Denote by \(i: A \to B\) the canonical injection. We then have \(i_*(\mathcal{F}) = (\mathcal{F}'(U), f_{\mathrm{UV}}')\) where for every open set U of B, \(\mathcal{F}'(U) = \mathcal{F}(U \cap A)\) , and for every pair (U, V) of open sets of B with \(U \subset V\) , \(f_{\mathrm{UV}}' = f_{(\mathrm{U} \cap A)(\mathrm{V} \cap A)}\) .

Now let \(\mathcal{G}\) be a presheaf on B. We define the inverse image of the presheaf \(\mathcal{G}\) under \(u\) , and denote by \(u^{*}(\mathcal{G})\) the sheaf \(\underline{\mathcal{C}}_{\mathrm{B}}(\mathrm{A};\mathrm{E}_{\mathcal{G}})\) on A of B-morphisms with values in the B-space \(\mathrm{E}_{\mathcal{G}}\) (I, p.~45, example 4). It is canonically isomorphic to the sheaf on A of sections of the A-étale space \(\mathrm{A} \times_{\mathrm{B}} \mathrm{E}_{\mathcal{G}}\) (I, p.~9, prop. 3). From this we deduce (I, p.~52, example 2) a canonical isomorphism \(\varphi\) of the étale space associated with \(u^{*}(\mathcal{G})\) onto the A-space \(\mathrm{A} \times_{\mathrm{B}} \mathrm{E}_{\mathcal{G}}\) . If moreover \(\mathcal{G}\) is a sheaf, then for every point \(a\) of A there is a canonical bijection \(\psi_{a}: u^{*}(\mathcal{G})_{a} \to \mathcal{G}_{u(a)}\) : for every open neighborhood U of \(a\) in A and every B-morphism \(f: \mathrm{U} \to \mathrm{E}_{\mathcal{G}}\) , we have

\[
\varphi ([ \mathrm{U}, f, a ]) = (a, f (a)), \quad \psi_ {a} ([ \mathrm{U}, f, a ]) = f (a).
\]

Let \(\mathcal{G}_1\) and \(\mathcal{G}_2\) be presheaves on B and \(\varphi\colon\mathcal{G}_1\to\mathcal{G}_2\) a morphism of presheaves. By base change of the morphism of B-étale spaces \(\mathrm{E}(\varphi)\colon\mathrm{E}_{\mathcal{G}_1}\to\mathrm{E}_{\mathcal{G}_2}\) we deduce a morphism of A-étale spaces \(\mathrm{A}\times_{\mathrm{B}}\mathrm{E}_{\mathcal{G}_1}\to\mathrm{A}\times_{\mathrm{B}}\mathrm{E}_{\mathcal{G}_2}\) , whence a morphism of sheaves on A, \(u^{*}(\varphi)\colon u^{*}(\mathcal{G}_{1})\to u^{*}(\mathcal{G}_{2})\) . Let \(\mathcal{G}_3\) be a presheaf on B and \(\psi\colon\mathcal{G}_2\to\mathcal{G}_3\) a morphism of presheaves. We then have the equality \(u^{*}(\psi\circ\varphi)=u^{*}(\psi)\circ u^{*}(\varphi)\) .

Let C be a topological space and let \(v\colon B \to C\) be a continuous map. If \(\mathcal{G}\) is a presheaf on C, the sheaves \(u^{*}(v^{*}(\mathcal{G}))\) and \((v \circ u)^{*}(\mathcal{G})\) are canonically identified (I, p.~5). If \(\varphi\colon\mathcal{G}_{1} \to \mathcal{G}_{2}\) is a morphism of presheaves on C, we moreover have \(u^{*}(v^{*}(\varphi)) = (v \circ u)^{*}(\varphi)\) .

Remark. --- The canonical morphism \(\sigma_{\mathcal{G}}\colon\mathcal{G}\to\widetilde{\mathcal{G}}\) of I, p.~53 corresponds in terms of étale spaces to the isomorphism \(j_{\mathcal{G}}\) of Proposition 2 of I, p.~53. It follows that \(u^{*}(\sigma_{\mathcal{G}})\) is an isomorphism. In particular, if \(A = B\) and \(u = \mathrm{Id}_A\) , the presheaf \(u^{*}(\mathcal{F})\) is the sheaf \(\widetilde{\mathcal{F}}\) .

Example 2. --- Let B be a topological space and A a subspace of B. Denote by \(i\colon\mathrm{A}\to\mathrm{B}\) the canonical injection. For every sheaf \(\mathcal{G}\) on B, denote by \(\mathcal{G}_{\mathrm{A}}\) the sheaf \(i^{*}(\mathcal{G})\) and call \(\mathcal{G}_{\mathrm{A}}\) the sheaf on A induced by the sheaf \(\mathcal{G}\) . The sheaf \(\mathcal{G}_{\mathrm{A}}\) is identified with the sheaf \(\underline{\mathcal{L}}(\mathrm{A};(\mathrm{E}_{\mathcal{G}})_{\mathrm{A}})\) of sections of the A-étale space induced by \(\mathrm{E}_{\mathcal{G}}\) over A.

Suppose that A is an open subspace of B, and let \(\mathcal{G}\) be a sheaf on B. By definition, the sheaf \(\mathcal{G}_{\mathrm{A}}\) is the sheaf \(\widetilde{\mathcal{G}}\mid \mathrm{A}\) deduced from \(\widetilde{\mathcal{G}}\) by restriction to the open set A (I, p.~43). Let \(\sigma_{\mathcal{G}}\colon \mathcal{G}\to \widetilde{\mathcal{G}}\) be the canonical isomorphism (I, p.~55, cor. 1). Then \(\sigma_{\mathcal{G}}\mid \mathrm{A}\) is a canonical isomorphism from the sheaf \(\mathcal{G}\mid \mathrm{A}\) onto the sheaf \(\mathcal{G}_{\mathrm{A}}\), called the canonical isomorphism of \(\mathcal{G}\mid \mathrm{A}\) onto \(\mathcal{G}_{\mathrm{A}}\) .

